% preamble_dsa.tex
\usepackage{fontspec}
\usepackage[spanish, es-noquoting, es-noshorthands]{babel}
\newfontfamily{\emojifont}{Segoe UI Emoji}[Renderer=HarfBuzz, BoldFont=*]
% Monospace con soporte de box-drawing (─│┼═) y flechas para el arte ASCII
\setmonofont{Consolas}[Scale=0.92]
% Glifos de caja/flechas/símbolos que la fuente roman no tiene:
% se remapean a Consolas cuando aparecen en texto normal (prosa, resumen)
\newfontfamily{\boxfont}{Consolas}
\usepackage{newunicodechar}
\newunicodechar{─}{{\boxfont ─}}
\newunicodechar{│}{{\boxfont │}}
\newunicodechar{┼}{{\boxfont ┼}}
\newunicodechar{┌}{{\boxfont ┌}}
\newunicodechar{┐}{{\boxfont ┐}}
\newunicodechar{└}{{\boxfont └}}
\newunicodechar{┘}{{\boxfont ┘}}
\newunicodechar{┴}{{\boxfont ┴}}
\newunicodechar{┬}{{\boxfont ┬}}
\newunicodechar{├}{{\boxfont ├}}
\newunicodechar{┤}{{\boxfont ┤}}
\newunicodechar{═}{{\boxfont ═}}
\newunicodechar{║}{{\boxfont ║}}
\newunicodechar{╱}{{\boxfont ╱}}
\newunicodechar{╲}{{\boxfont ╲}}
\newunicodechar{→}{{\boxfont →}}
\newunicodechar{←}{{\boxfont ←}}
\newunicodechar{↑}{{\boxfont ↑}}
\newunicodechar{↓}{{\boxfont ↓}}
\newunicodechar{↔}{{\boxfont ↔}}
\newunicodechar{⇄}{{\boxfont ⇄}}
\newunicodechar{◄}{{\boxfont ◄}}
\newunicodechar{►}{{\boxfont ►}}
\newunicodechar{▲}{{\boxfont ▲}}
\newunicodechar{▼}{{\boxfont ▼}}
\newunicodechar{✓}{{\boxfont ✓}}
\newunicodechar{✗}{{\boxfont ✗}}
\newunicodechar{≤}{{\boxfont ≤}}
\newunicodechar{≥}{{\boxfont ≥}}
\newunicodechar{≠}{{\boxfont ≠}}
\newunicodechar{▶}{{\boxfont ▶}}
\newunicodechar{◀}{{\boxfont ◀}}
\newcommand{\emj}[1]{{\emojifont #1}}
\usepackage[a4paper, top=2.5cm, bottom=2.5cm, left=3cm, right=2.5cm]{geometry}
\usepackage{fancyhdr}
\usepackage{titlesec}
\usepackage{hyperref}
\usepackage{xcolor}
\usepackage{listings}
\usepackage{tcolorbox}
\usepackage{graphicx}
\usepackage{booktabs}
\usepackage{enumitem}
\usepackage{parskip}
\raggedbottom

\tcbuselibrary{skins, breakable, listings}

% --- Colores Neon ---
\definecolor{n-cyan}{RGB}{0, 255, 255}
\definecolor{n-lime}{RGB}{0, 255, 0}
\definecolor{n-pink}{RGB}{255, 0, 255}
\definecolor{n-yellow}{RGB}{255, 255, 0}
\definecolor{n-orange}{RGB}{255, 128, 0}

% --- Colores de Apoyo ---
\definecolor{azuloscu}{RGB}{30, 64, 175}
\definecolor{verdeoscu}{RGB}{6, 95, 70}
\definecolor{naranjaoscu}{RGB}{154, 52, 18}
\definecolor{rojooscu}{RGB}{153, 27, 27}
\definecolor{amarillooscu}{RGB}{113, 63, 18}
\definecolor{moradooscu}{RGB}{88, 28, 135}
\definecolor{cianoscuro}{RGB}{14, 116, 144}
\definecolor{grisclaro}{RGB}{200, 200, 200}
\definecolor{codegray}{RGB}{130, 130, 130}

% --- MODO OSCURO GLOBAL ---
\definecolor{fondogris}{RGB}{30, 30, 35}
\pagecolor{fondogris}
\color{white}
\makeatletter
\colorlet{normalcolor@tmp}{white}
\def\normalcolor{\color{normalcolor@tmp}}
\makeatother

\hypersetup{
  colorlinks=true,
  linkcolor=n-cyan,
  urlcolor=n-lime,
  citecolor=grisclaro,
  pdftitle={Apuntes DSA - Data Structures \& Algorithms},
  pdfauthor={Aldo Zepeda Montoya},
}
% Prevent emoji commands from breaking PDF bookmarks
\pdfstringdefDisableCommands{%
  \def\emj#1{}%
}

% --- Estilo de Código C++ (paleta VS Code Dark+) ---
\definecolor{vscKeyword}{RGB}{86, 156, 214}    % azul: int, double, const, class, using...
\definecolor{vscControl}{RGB}{197, 134, 192}   % púrpura: if, for, return, try, catch + pseudocódigo
\definecolor{vscType}{RGB}{78, 201, 176}       % teal: std, vector, string, map (clases/STL)
\definecolor{vscString}{RGB}{206, 145, 120}    % salmón: strings
\definecolor{vscComment}{RGB}{106, 153, 85}    % verde: comentarios
\lstdefinestyle{cpp}{
  language=C++,
  basicstyle=\ttfamily\small\color{white},    % operadores y variables: blanco (como VS Code)
  % Clase 1: keywords de tipo/declaración (int, double, const, class, using...) -> azul
  keywordstyle=\color{vscKeyword},
  % Clase 2: clases y contenedores STL -> teal
  keywords=[2]{std,string,vector,queue,deque,stack,priority_queue,unordered_map,unordered_set,tuple,size_t,greater,less,iterator},
  keywordstyle=[2]\color{vscType},
  % Clase 3: control de flujo (C++ y pseudocódigo español) -> púrpura
  keywords=[3]{if,else,for,while,do,return,switch,case,break,continue,goto,try,catch,throw,SI,SINO,MIENTRAS,PARA,REPETIR,HASTA,DESDE,funcion,romper,continuar,retornar,devolver,cada,ENTONCES,VECES},
  keywordstyle=[3]\color{vscControl},
  % Clase 4: constantes/macros -> azul (como VS Code)
  keywords=[4]{NULL,INFINITO,VERDADERO,FALSO,INF},
  keywordstyle=[4]\color{vscKeyword},
  % Preprocesador (#include, #define) -> púrpura
  moredelim=[l][\color{vscControl}]{\#},
  commentstyle=\color{vscComment},
  stringstyle=\color{vscString},
  numberstyle=\tiny\color{codegray},
  showstringspaces=false,
  breaklines=true,
  frame=single,
  framerule=0.6pt,
  rulecolor=\color{n-cyan},
  backgroundcolor=\color{black!94!white},
  numbers=left,
  numbersep=8pt,
  tabsize=2,
  inputencoding=utf8,
  extendedchars=true,
  literate=
    {á}{{\'{a}}}1 {é}{{\'{e}}}1 {í}{{\'{i}}}1
    {ó}{{\'{o}}}1 {ú}{{\'{u}}}1 {ñ}{{\~{n}}}1
    {Á}{{\'{A}}}1 {É}{{\'{E}}}1 {Í}{{\'{I}}}1
    {Ó}{{\'{O}}}1 {Ú}{{\'{U}}}1 {Ñ}{{\~{N}}}1
    {¿}{{\textquestiondown}}1 {¡}{{\textexclamdown}}1
}
\lstset{style=cpp}

% --- tcolorbox environments ---
\newtcolorbox{dsaquees}{
  colback=azuloscu!25!black,
  colframe=n-cyan,
  colupper=white,
  fonttitle=\bfseries\large,
  title={\emj{🤔}~¿Qué es?},
  breakable,
  before upper={\parskip=4pt}
}

\newtcolorbox{dsanino}{
  colback=moradooscu!20!black,
  colframe=n-lime,
  colupper=white,
  fonttitle=\bfseries\large,
  title={\emj{🧒}~Explicado para un niño de 12},
  breakable,
  before upper={\parskip=4pt}
}

\newtcolorbox{dsaotraforma}{
  colback=moradooscu!25!black,
  colframe=n-pink,
  colupper=white,
  fonttitle=\bfseries\large,
  title={\emj{💬}~Dicho de otra forma},
  breakable,
  before upper={\parskip=4pt}
}

\newtcolorbox{dsadiagrama}{
  colback=verdeoscu!25!black,
  colframe=n-lime,
  colupper=white,
  fonttitle=\bfseries\large,
  title={\emj{🎨}~Diagrama},
  breakable,
  before upper={\parskip=4pt}
}

\newtcolorbox{dsacomofunciona}{
  colback=cianoscuro!25!black,
  colframe=n-cyan,
  colupper=white,
  fonttitle=\bfseries\large,
  title={\emj{⚙️}~¿Cómo funciona?},
  breakable,
  before upper={\parskip=4pt}
}

\newtcolorbox{dsacomplejidad}{
  colback=amarillooscu!25!black,
  colframe=n-yellow,
  colupper=white,
  fonttitle=\bfseries\large,
  title={\emj{📊}~Complejidad Big-O},
  breakable,
  before upper={\parskip=4pt}
}

\newtcolorbox{dsaentrevistas}{
  colback=naranjaoscu!25!black,
  colframe=n-orange,
  colupper=white,
  fonttitle=\bfseries\large,
  title={\emj{🎯}~En entrevistas},
  breakable,
  before upper={\parskip=4pt}
}

\newtcolorbox{dsaresumen}{
  colback=rojooscu!25!black,
  colframe=n-pink,
  colupper=white,
  fonttitle=\bfseries\large,
  title={\emj{📋}~Resumen rápido},
  breakable,
  before upper={\parskip=4pt}
}

% --- Formato Capítulos y Partes ---
\titleformat{\part}[display]
  {\normalfont\Huge\bfseries\color{n-yellow}\centering}{}{0pt}{\Huge}
\titleformat{\chapter}[display]
  {\normalfont\huge\bfseries\color{n-cyan}}{\chaptertitlename\ \thechapter}{12pt}{\Huge}
\titlespacing*{\chapter}{0pt}{0pt}{16pt}
\titleformat{\section}
  {\normalfont\large\bfseries\color{n-lime}}{}{0pt}{}

% --- Encabezados ---
\pagestyle{fancy}
\fancyhf{}
\fancyhead[LE,RO]{\color{n-cyan}\thepage}
\fancyhead[LO]{\color{grisclaro}\nouppercase{\rightmark}}
\fancyhead[RE]{\color{grisclaro}\nouppercase{\leftmark}}
\renewcommand{\headrulewidth}{0.4pt}
\renewcommand{\headrule}{\hbox to\headwidth{\color{n-cyan}\leaders\hrule height \headrulewidth\hfill}}
